\section{Additional assessment of the physics revision}
\label{sec:physics-lean-assessment}

The new physics sections provide much more explicit mathematical data:
character coefficients, Virasoro-primary multiplicities, spin-resolved
spectra, thermal intervals, low-weight couplings, rational correlators,
fusion coefficients, defects, and boundary partition functions. None
of the Python producers in \artifact{c8_physics_20260927/} is a Lean
proof or has a Lean-verified implementation. Their exact arithmetic
and internal checks improve reviewability and supply concrete inputs
for formalization; they do not enlarge the scopes of the five existing
CFT declarations listed in the preceding assessment. A separate new
file adds the two small thermal-comparison lemmas stated below, bringing
the selected CFT coverage to seven principal declarations in three
files. The five baseline declarations retain their original scope.

\subsection{Current manuscript anchors and historical concordance}

The linked numbers, titles, and pages in the following table refer to
the current manuscript. The first column is only a concordance for
readers of the archived 55-page PDF; its former section numbers are
not operative references in this guide. The current document places
the mathematical development in appendices after the physics sections.

\begin{longtable}{@{}P{0.09\linewidth}P{0.59\linewidth}P{0.24\linewidth}@{}}
\toprule
Archived label & Current manuscript title & Current anchor\\
\midrule\endfirsthead
\toprule Archived label & Current manuscript title & Current anchor\\
\midrule\endhead\bottomrule\endfoot
1 & \msnameref{sec:statement}
& Section~\msref{sec:statement}, p.~\mspageref{sec:statement}\\
2 & \msnameref{sec:construction}
& Appendix~\msref{sec:construction}, p.~\mspageref{sec:construction}\\
3 & \msnameref{sec:finite}
& Appendix~\msref{sec:finite}, p.~\mspageref{sec:finite}\\
4 & \msnameref{sec:gate4}
& Appendix~\msref{sec:gate4}, p.~\mspageref{sec:gate4}\\
5 & \msnameref{sec:gate5}
& Appendix~\msref{sec:gate5}, p.~\mspageref{sec:gate5}\\
6 & \msnameref{sec:properness}
& Appendix~\msref{sec:properness}, p.~\mspageref{sec:properness}\\
7 & \msnameref{sec:regularity}
& Appendix~\msref{sec:regularity}, p.~\mspageref{sec:regularity}\\
8 & \msnameref{sec:branching}
& Appendix~\msref{sec:branching}, p.~\mspageref{sec:branching}\\
9 & \msnameref{sec:characters}
& Appendix~\msref{sec:characters}, p.~\mspageref{sec:characters}\\
10 & \msnameref{sec:analytic-unitarity}
& Appendix~\msref{sec:analytic-unitarity}, p.~\mspageref{sec:analytic-unitarity}\\
11 & \msnameref{sec:analytic-category}
& Appendix~\msref{sec:analytic-category}, p.~\mspageref{sec:analytic-category}\\
12 & \msnameref{sec:analytic-full-theory}
& Appendix~\msref{sec:analytic-full-theory}, p.~\mspageref{sec:analytic-full-theory}\\
13 & \msnameref{sec:comparison}
& Appendix~\msref{sec:comparison}, p.~\mspageref{sec:comparison}\\
A & \msnameref{sec:ledger}
& Appendix~\msref{sec:ledger}, p.~\mspageref{sec:ledger}\\
B & \msnameref{sec:formalization}
& Appendix~\msref{sec:formalization}, p.~\mspageref{sec:formalization}\\
C & \msnameref{sec:compilation}
& Appendix~\msref{sec:compilation}, p.~\mspageref{sec:compilation}\\
\end{longtable}

The added physics material has the following current anchors.
``Focused'', ``substantial'', and ``foundational'' retain the
dependency-based meanings stated earlier.

\begin{longtable}{@{}P{0.70\linewidth}P{0.21\linewidth}@{}}
\toprule
Current section title & Current anchor\\
\midrule\endfirsthead
\toprule Current section title & Current anchor\\
\midrule\endhead\bottomrule\endfoot
\msnameref{sec:physics-spectrum}
& Section~\msref{sec:physics-spectrum}, p.~\mspageref{sec:physics-spectrum}\\
\msnameref{sec:physics-correlators}
& Section~\msref{sec:physics-correlators}, p.~\mspageref{sec:physics-correlators}\\
\msnameref{sec:physics-couplings}
& Section~\msref{sec:physics-couplings}, p.~\mspageref{sec:physics-couplings}\\
\msnameref{sec:physics-defects}
& Section~\msref{sec:physics-defects}, p.~\mspageref{sec:physics-defects}\\
\msnameref{sec:physics-H3}
& Section~\msref{sec:physics-H3}, p.~\mspageref{sec:physics-H3}\\
\msnameref{sec:physics-data-guide}
& Section~\msref{sec:physics-data-guide}, p.~\mspageref{sec:physics-data-guide}\\
\end{longtable}

\subsection{Characters, Virasoro primaries, and the full spectrum}

Section~\msref{sec:physics-spectrum},
\emph{\msnameref{sec:physics-spectrum}}, is in
\artifact{c8_physics_20260927/spectrum.tex}. Its producer
\artifact{spectrum.py} checks the published initial coefficients,
the rational matrix recurrence, nonnegative integrality through the
requested grade, descendant reconstruction, parity of the resulting
spin spectrum, and an independent $E_4/\eta^8$ identity. The default
data include 25 coefficients for each of twelve characters, complete
full-CFT spectra through $\Delta=24$, and an independent $E_8$ spectrum.
These are finite arithmetic conclusions conditional on the input
matrices and their character interpretation.

\begin{longtable}{@{}P{0.26\linewidth}P{0.64\linewidth}@{}}
\toprule
Physics argument & Missing Lean theorem and difficulty\\
\midrule\endfirsthead
\toprule Physics argument & Missing Lean theorem and difficulty\\
\midrule\endhead\bottomrule\endfoot
Matrix recurrence
& Define the twelve-dimensional exponent/characteristic data and a
finite rational recurrence, prove nonzero denominators, and prove that
its finite output agrees with the formal differential equation.
This is focused to substantial. Identifying the resulting vector with
the actual VOA trace functions still requires the foundational
character theorem assessed under Appendix~\msref{sec:characters}.\\
Virasoro-primary subtraction
& The finite Euler-product convolution and its inverse are focused.
Their interpretation needs the irreducible $c=8$ Virasoro characters,
the Kac determinant or an equivalent irreducibility theorem, and the
unitary highest-weight decomposition of each module. These are
substantial to foundational representation-theoretic inputs.\\
Complete spin-resolved cutoff
& Given nonnegative lowest weights and coefficient functions, prove
the product formula and that no grade above $N$ contributes to
$\Delta\leq N$. The finite sums, parity identity, and relevant/marginal
counts are focused. The tensor-product Hilbert-space interpretation
and the diagonal modular invariant are additional premises.\\
$E_8$ comparison and invariant counts
& A finite identity between the oscillator series and the selected
character combination is focused. It does not identify the extension
algebra. Counts $1488=6\cdot248$ for one chiral Haagerup action and
$36=6^2$ for both independent actions additionally require the precise
invariant-space theorem. That categorical interpretation belongs to
the separate defect formalization below.\\
\end{longtable}

There is useful existing power-series infrastructure.
\leanid{PowerSeries.coeff_mk} and \leanid{PowerSeries.coeff_mul} occur
in \artifact{Mathlib/RingTheory/PowerSeries/Basic.lean}.
The file
\artifact{Mathlib/Combinatorics/Enumerative/Pentagonal/PowerSeries.lean}
contains the Euler-product development, including
\leanid{PowerSeries.coeff_prod_one_sub_X_pow_eventually_eq} and
\leanid{PowerSeries.WithPiTopology.hasProd_one_sub_X_pow}.
These declarations were read in the pinned checkout. They support
formal coefficient manipulations; they do not state that these
coefficients count this CFT's Virasoro primaries.

An economical formalization would first take an abstract sequence of
graded dimensions and prove the descendant-removal and tensor-product
counting lemmas. A second layer would connect that sequence to the
Evans--Gannon recurrence. The final interpretation theorem must connect
both layers to the actual modules. Merely importing the printed
integers as definitions and proving that their sum is $2941$ would
check arithmetic but leave the substantial spectral assertion open.

\subsection{Thermal bounds and their numerical intervals}

Proposition~\msref{prop:physics-modular-tail}
(p.~\mspageref{prop:physics-modular-tail}) proves the modular heavy-tail
estimate in Section~\msref{sec:physics-spectrum}. For
$E=\Delta-2/3$, an exact cutoff $D>2/3$, and $b'=4\pi^2/b$, it uses
\[
 Z=L_D+H_D,\qquad Z(b)=Z(b'),\qquad
 H_D(b)\leq e^{-(b-b')(D-2/3)}H_D(b')
\]
to bound the tail at $b>2\pi$. Two real-number steps are now checked in
\artifact{lean/HaagerupIzumi/CFT/PhysicsBounds.lean}, in namespace
\leanid{HaagerupIzumi.CFT}:
\begin{itemize}
\item \leanid{modular_heavy_tail_comparison} assumes
$L+H=L_d+H_d$, $H\leq rH_d$, and $0\leq r<1$, and proves
$H\leq r(L-L_d)/(1-r)$;
\item \leanid{selfdual_tail_first_moment_comparison} assumes
$E_{\rm cut}>0$ and $E_{\rm cut}H\leq L'$, and proves
$H\leq L'/E_{\rm cut}$.
\end{itemize}
The successful guarded compiler record is
\artifact{c8_physics_20260927/physics_lean.log}, with sampled peak
group RSS $1.549$~GiB. Both axiom reports contain only
\leanid{propext}, \leanid{Classical.choice}, and \leanid{Quot.sound};
the source contains neither \texttt{sorry} nor new axioms. The
nonnegative contraction range is retained explicitly, although the
first rearrangement only needs $r<1$ once its comparison is supplied.

These lemmas do not derive the comparison from a spectral sum or
identify $r$ with the exponential in the application. The actual
infinite-sum theorem remains more substantial: it needs convergence,
termwise positivity, completeness of the finite cutoff, and the
modular functional equation. Its self-dual form also needs
differentiability and the identity $Z'(2\pi)=0$. The first and second
moment bounds require justified moment convergence and the elementary
optimization of $E^ke^{-aE}$ on a positive half-line.

The Python output is a conditional interval certificate. Its positive
spectrum and modular invariance are mathematical premises, not facts
established by evaluating finitely many coefficients. The implementation
uses directed decimal arithmetic, a fixed rational enclosure of $\pi$,
and correctly rounded exponential/logarithm functions; the displayed
six-decimal table is accepted only when every rigorous interval lies
within $10^{-6}$ of its printed value. None of that implementation is
currently kernel checked.

For a Lean version, a small rational interval checker with proved
exponential and logarithm enclosures would be preferable to treating
Python's decimal library as an axiom. One can separately formalize the
analytic tail theorem over the real numbers and feed it certified
finite moment bounds. Existing real analysis and infinite-sum
infrastructure is relevant, but the complete pipeline from characters
to the saved numerical intervals is new work. The two checked
implications are the completed local algebraic part; sound interval
evaluation and its analytic interpretation
are substantial. The CFT assumptions still inherit the foundational
requirements of the earlier sections.

\subsection{Currents, lattice correlators, and low-weight couplings}

The correlator and coupling material in
Sections~\msref{sec:physics-correlators} and
\msref{sec:physics-couplings} is supplied by
\artifact{c8_physics_20260927/correlators.tex} and
\artifact{c8_physics_20260927/couplings.tex}. Its data distinguish
direct finite lattice-mode products, rational identities reconstructed
from saved singular OPEs, and coefficients completed using current
Ward identities. Those three forms of evidence have different
interpretation obligations; the final reports specify the supplied
bases and direct probes.

In particular the completed weight-two dataset in
\artifact{c8_physics_20260927/full_weight_two/} uses all $113$ chiral
Virasoro primaries. Its saved report records $1{,}442{,}897$ ordered
tensor slots, $18{,}365$ nonzero entries, $51{,}076$ current Ward
checks, and $33$ additional direct lattice-mode probes beyond the
seed calculation. This supplies substantially more explicit input
than a list of summed OPE coefficients. Those successful Python
checks remain separate from the Lean obligations below.

The following finite targets are realistic once the coefficient
algebra and basis conventions are represented in Lean:
\begin{itemize}
\item singular current brackets, levels, gauge factors, pairings, and
the orthogonality and positivity of the stated finite Gram matrices;
\item polynomial identities obtained by clearing denominators in
the displayed rational four-point functions, with collision loci
excluded when interpreting them as functions;
\item exact extraction of three-point coefficients and their squared
unit-norm values, retaining the specified invariant/Hermitian pairing
and adjoint conventions;
\item the finite Ward recursion, its compatibility with both current
factors, and uniqueness of its completion from the supplied seed slice;
\item the completeness of the chosen low-weight primary subspace,
provided its exact rank and descendant-complement certificates.
\end{itemize}
These are substantial finite algebra rather than a need to formalize
all CFT at once. The bridge from encoded lattice words to actual mode
operators is the same missing bridge described for the finite gates.
A successful check of a completed tensor against Ward identities
establishes consistency of the completion. Its identification with
all actual mode products also needs a proved determination theorem
and authenticated seed products; a finite selection of direct probes
does not by itself check every omitted product.

Beyond that finite layer, one needs the state-field correspondence,
locality, conformal covariance, OPE convergence or formal expansion
identities, and an explanation of how the chiral functions assemble
into local bulk correlators. The known Potts and parafermion inputs
also need their formal embeddings. These are foundational. A chiral
vacuum correlator or a chosen lattice intertwiner block is not by
itself the four-point function of the lightest diagonal bulk scalar.
The formal statement must preserve the external fields and their
normalization rather than silently promote one kind of block to the
other.

No claim that a dimension-two scalar is exactly marginal follows
from its existence or its three-point coefficient alone. A formal
statement about a conformal manifold would additionally require a
precise deformation theory, beta functions, and the appropriate
all-orders argument. That subject is outside both the present Lean
files and the finite arithmetic producer.

\subsection{Modular data, defects, and conformal boundaries}

Sections~\msref{sec:physics-defects} and \msref{sec:physics-H3},
\emph{\msnameref{sec:physics-defects}} and
\emph{\msnameref{sec:physics-H3}}, are in
\artifact{c8_physics_20260927/defects.tex}. Their
producer \artifact{defects.py} computes in
$\mathbb Q(\zeta_{13})$, with exact Verlinde coefficients, modular
relations, the forgetful multiplicities, and boundary/defect spectral
tables. These checks establish finite algebraic properties of the
input matrices and based rings, not the existence of the required
associator or a conformal defect.

\begin{longtable}{@{}P{0.27\linewidth}P{0.63\linewidth}@{}}
\toprule
Physics argument & Missing Lean theorem and difficulty\\
\midrule\endfirsthead
\toprule Physics argument & Missing Lean theorem and difficulty\\
\midrule\endhead\bottomrule\endfoot
Cyclotomic modular arithmetic
& Prove that the twelve-coordinate quotient representation is a
field embedded in $\mathbb C$, fix the primitive root and conjugation,
and verify the Gauss-sum square root and sign. Then $S^2=1$, the
twist relation, Verlinde integrality, and fusion associativity are
substantial but bounded exact computations.\\
Diagonal Verlinde lines
& Given a unitary modular tensor category and the diagonal full CFT,
derive line eigenvalues, twisted character matrices, and their
sector multiplicities. The scalar sums are focused; the theorem that
these matrices describe actual topological defects requires the
categorical/full-CFT construction and is foundational.\\
Cardy boundaries and annuli
& Formalize the boundary-state normalization, open/closed-channel
identity, and the positive square roots entering $g$ factors.
Finite character multiplicities are focused. Their interpretation
as boundary Hilbert spaces requires the Cardy/extension theorem.
Formal verification of annulus integers is not that theorem.\\
Selected chiral $\mathcal H_3$ action
& Prove the center-forgetful correspondence for the chosen extension,
the bimodule/full-defect equivalence, and the identification of tube
multiplicity spaces with $M_\mu\otimes\overline{H'}$.
These are foundational categorical inputs. Once supplied, the
$6\times12$ forgetful ring-map check and charge matrices are
substantial finite algebra.\\
Invariant and charged gaps
& Under the correct action, derive the fixed spaces and count their
lowest scalar grades. This is focused after the sector assignment.
It must distinguish one chiral action from both independent factors,
and tube multiplicity spaces from their multiple copies in actual
defect Hilbert spaces. Modular covariance alone is insufficient.\\
Six extension boundaries
& Verify the finite annulus coefficients from the forgetful functor
and the six-object fusion ring. The data-to-category interface must
retain the actual module category, duals, and boundary normalization.
The finite arithmetic is substantial; constructing these boundaries
as part of the full CFT remains foundational.\\
\end{longtable}

The pinned checkout defines \leanid{CyclotomicField} and
\leanid{IsCyclotomicExtension} in
\artifact{Mathlib/NumberTheory/Cyclotomic/Basic.lean}; it also has
primitive-root, complex-root-of-unity, and cyclotomic-polynomial
developments. Thus the field-theory layer need not start from scratch.
It still takes work to prove that the producer's concrete coordinate
arithmetic and chosen complex embedding agree with those structures.
An arbitrary root of the polynomial without the intended complex
embedding does not determine the positive quantum dimensions used
by the physics interpretation.

\subsection{Bootstrap placement and observable scope}

The physics reading guide in Section~\msref{sec:physics-data-guide}
distinguishes the diagonal full theory from
the selected holomorphic extension and states which published bootstrap
plots use which symmetry sectors. Formalizing the label correspondence,
the central-charge convention $c_L=c_R=8$, and a rational gap extracted
from a proved spectrum would be focused. It would not formally verify
a published numerical exclusion curve. That requires the original
functional, its positivity on every required conformal block, the
truncation/error bounds, and a sound numerical or exact certificate.
A plotted point read from a figure is evidence for comparison, not a
kernel-checked exclusion or inclusion theorem.

The observable inventory is also an interface specification. Characters
determine unflavoured multiplicities and, with analytic control,
partition functions. They do not determine every normalized OPE
coefficient or defect eigenvector. The new coupling and defect datasets
close specific parts of that gap; each needs its own sound checker and
interpretation theorem. Keeping those claims separate is essential to
an eventual formal treatment that a physicist can use without reading
the implementation.

\subsection{A concrete order of formalization}

A useful first deliverable is a small exact coefficient library covering
the rational recurrence, finite Euler product, primary reconstruction,
and cutoff-completeness lemma. Its theorem should expose the input
character interpretation explicitly. A second deliverable is a sound
finite algebra checker for the cyclotomic modular data and the supplied
low-weight tensors, with an explicit interpretation map. A third is
the real-analytic tail theorem and a certified interval evaluator.
Each can be independently reviewed before the foundational VOA and
defect interfaces are complete.

Completion criteria must be semantic: the formal theorem should refer
to the actual character coefficients, mode products, or represented
sectors through proved maps. File hashes, successful Python assertions,
or a Lean theorem about unrelated constants cannot replace that
connection. The two explicitly stated thermal-comparison lemmas are
the additional Lean results of this physics revision; the much larger
analytic and physical interpretations remain to be formalized.
